\documentclass{practice}

\title{6}
\date{\today}

\usepackage{crysymb}

\usepackage{fancyvrb}
\fvset{listparameters=\setlength{\topsep}{0pt}\setlength{\partopsep}{0pt}}

\usepackage{listings}

\lstset{
basicstyle=\small\ttfamily,
columns=flexible,
breaklines=true
}

\begin{document}
\maketitle

\begin{task}{OpenSSL}
  \textit{Preface.}
  For working with public-key cryptography keypairs, OpenSSL provides the following tools:
  \begin{itemize}
    \item \href{https://docs.openssl.org/master/man1/openssl-genpkey/}{\texttt{genpkey}}: utility for generating a private key (RSA, EC) or parameters (DH, DSA, EC)
    \begin{itemize}
      \item Alternatively: \href{https://docs.openssl.org/master/man1/openssl-genrsa/}{\texttt{genrsa}}, \href{https://docs.openssl.org/master/man1/openssl-dhparam/}{\texttt{dhparam}}, \href{https://docs.openssl.org/master/man1/openssl-ecparam/}{\texttt{ecparam}}
    \end{itemize}
    \item \href{https://docs.openssl.org/master/man1/openssl-pkey/}{\texttt{pkey}}: utility for processing keys, e.g. change their format (PEM, DER) or passphrase
    \begin{itemize}
      \item Alternatively: \href{https://docs.openssl.org/master/man1/openssl-rsa/}{\texttt{rsa}}, \href{https://docs.openssl.org/master/man1/openssl-ec}{\texttt{ec}}
    \end{itemize}
    \item \href{https://docs.openssl.org/master/man1/openssl-pkeyutl/}{\texttt{pkeyutl}}: utility for actual operations with keys (e.g. encryption or signing)
    \begin{itemize}
      \item \href{https://docs.openssl.org/master/man1/openssl-rsautl/}{\texttt{rsautl}} is deprecated, do not use it
    \end{itemize}
  \end{itemize}

  To generate parameters instead of a private key with the \texttt{genpkey} utility, you must use the \texttt{-genparam} option.
  For example, if you one day feel the strong urge to generate a large prime before going to bed, you can use the following command:
  \begin{Verbatim}
openssl genpkey -genparam -algorithm dh \
    -pkeyopt dh_paramgen_prime_len:3072
  \end{Verbatim}

  If you wonder what the dots, plus signs, and asterisks that appear mean, then for DH there is a nice InfoSec Stack Exchange \href{https://security.stackexchange.com/a/140639}{\textit{post}} on it.

  In practice, perhaps the primary use case for generating parameters in this manner with OpenSSL is to publish the parameters or share them with other entities.
  These parameters then form the basis of subsequent procedures\footnotemark{}.
  For example, to generate a private key with OpenSSL based on a set of parameters, you can use the \texttt{-paramfile} option.
  \footnotetext{E.g. the frowned upon approach of generating custom Diffie-Hellman parameters for TLS 1.2.}

  \textit{Task.}
  For this task, familiarise yourself with the OpenSSL documentation.
  Then, generate a password-protected P-384 key encrypted with AES-128 in CBC mode.\footnotemark{}
  \footnotetext{\texttt{genpkey} does not support AEAD for some reason, and AES in CTR mode has an implementation bug for this use.
  Marvels of OpenSSL\dots{}.}
\end{task}

\begin{task}{HPKE}
  \textit{Preface.}
  In the lecture we briefly discussed IES and ECIES.
  However, as you may recall, the various (!) ECIES specifications are not freely accessible, and IES itself does not have a single \enquote{canonical} standard.
  The modern alternative is \emph{Hybrid Public Key Encryption} (HPKE), specified in \href{https://datatracker.ietf.org/doc/html/rfc9180}{RFC 9180}.

  In its commonly used form, HPKE combines a KEM, a KDF, and an AEAD scheme.
  To encrypt a message, the sender must know the recipient's public key -- whether long-term or ephemeral (e.g. to achieve forward secrecy).

  You may then ask, should the sender be authenticated as well?
  HPKE supports four \emph{modes}, which differ in how (and whether) the sender is authenticated:
  \begin{itemize}
    \item Base -- The sender generates a fresh ephemeral keypair for each message and remains anonymous (at the scheme level) to the recipient.
    \item PSK -- \emph{Pre-shared key}.
    A symmetric key known to both parties is incorporated into the key \enquote{schedule}.
    This provides symmetric authenticity guarantees.
    \item Auth -- The sender additionally uses a static (long-term) key pair to implicitly authenticate themselves to the recipient.
    \item Auth + PSK -- Auth mode with an additional PSK mixed into the key schedule.
  \end{itemize}

  \textit{Task.}
  pyca/cryptography provides HPKE through the \href{https://cryptography.io/en/latest/hazmat/primitives/hpke/}{\texttt{hpke}} module.
  A ciphersuite is chosen by combining a KEM, a KDF, and an AEAD into a \texttt{Suite}, which then exposes single-shot \texttt{encrypt} and \texttt{decrypt} methods.
  Note that pyca/cryptography only implements the base mode.

  For this task, read through the package documentation and then do the following:
  \begin{enumerate}
    \item Generate your own P-384 keypair.\footnote{In practice, X25519 is the most widely deployed HPKE KEM.}
    \item Submit your public key to the Discord bot, and receive a ciphertext.
    \item Decrypt the AES-128-GCM encrypted ciphertext with the context (info)
    \begin{Verbatim}
ICS0036 week 6
    \end{Verbatim}
    The server is using HPKE in base mode with HKDF-SHA384 as the KDF.
  \end{enumerate}
\end{task}

\begin{task}{HPKE from scratch}
  \textit{Preface.}
  In the previous task, \texttt{Suite.encrypt} and \texttt{Suite.decrypt} were black boxes.
  For the sake of understanding, let us implement our own HPKE encryption in base mode.
  As in the server's ciphersuite, we use a Diffie-Hellman based KEM (DHKEM).

  \textit{Task.}
  Complete the \texttt{encrypt} function in the template.
  Each step below corresponds to a step in the template.
  Steps 1--3 are the DHKEM encapsulation (slide 38 of the lecture), and steps 4--6 are the generic HPKE steps (slide 37).
  \begin{enumerate}
    \item Compute $R \gets rG$ with $r \getsu \ZZ_n^*$.

    This is the sender's half of a Diffie-Hellman key exchange.
    The recipient's half is their public key $K_B$ which they publish in advance.
    $R$ is the \emph{encapsulation} $\mathit{enc}$ that the recipient later needs to recover the shared secret.

    \textit{Question:} if $K_B$ is a long-term key, what happens to past messages when $k_B$ leaks?

    \item Compute the shared point $S \gets rK_B$.

    \item Derive the shared secret $s \gets \KDF(S, R||K_B)$.

    $S$ is a curve point, not a uniformly random bit string, so it cannot be used as a key directly.
    
    Mixing in $R||K_B$ ties the shared secret to this exact exchange.
    Without it, an attacker could modify $R$ in a way that leads to the same key, producing a different ciphertext that still decrypts correctly (i.e. the ciphertext would be \emph{malleable}).

    \item Derive the AEAD key and nonce $k, \mathit{nonce} \gets \KDF(s, \mathit{info})$.

    \textit{Question:} why a second KDF step instead of deriving everything from $S$ in step 3?
    
    Hint: is step 3 part of every KEM?
    Compare slides 37 and 38.

    \item Encrypt $c \gets \ENC_k(\mathit{nonce}, \mathit{aad}, m)$ with AES-128-GCM.

    \item Output $\mathit{enc}||c$, i.e. $R||c$.

    \textit{Question:} why must $R$ be sent along with the ciphertext, but not the nonce?
  \end{enumerate}

  The template decrypts your ciphertext with pyca/cryptography.
  If your implementation is correct, you will see your original message.
\end{task}
\end{document}
